% !TEX program = xelatex
% Gerado por artigo/gerar_uso_ia.py a partir de artigo/declaracao-uso-ia.md e artigo/relatorio-uso-ia.md.
% Não edite à mão: altere os .md e rode `python artigo/gerar_uso_ia.py --pdf`.
\documentclass[12pt,a4paper]{article}
\usepackage[top=3cm,bottom=2cm,left=3cm,right=2cm,headsep=0.8cm]{geometry}
\usepackage{iftex}
\usepackage{amsmath}
\usepackage[brazilian]{babel}
\ifPDFTeX
  % pdfLaTeX: Times pelo newtx (ou pelo mathptmx, se o newtx não estiver instalado)
  \usepackage[T1]{fontenc}
  \usepackage[utf8]{inputenc}
  \IfFileExists{newtxtext.sty}{\usepackage{newtxtext,newtxmath}}{\usepackage{mathptmx}}
\else
  % XeLaTeX/LuaLaTeX: Times New Roman 12, como pedem as instruções da disciplina. Onde ela não
  % estiver instalada (Overleaf, Linux), usa a TeX Gyre Termes, clone métrico da Times do TeX Live.
  \usepackage{unicode-math}
  \IfFontExistsTF{Times New Roman}
    {\setmainfont{Times New Roman}}
    {\setmainfont{texgyretermes}[Extension=.otf, UprightFont=*-regular, BoldFont=*-bold,
       ItalicFont=*-italic, BoldItalicFont=*-bolditalic]}
  \setmathfont{texgyretermes-math.otf}
\fi
\usepackage{microtype}
\usepackage{indentfirst}  % recua também o primeiro parágrafo de cada seção (uso brasileiro)
\usepackage{graphicx}
\usepackage{calc}
\usepackage{array}
\usepackage{longtable}
\usepackage{booktabs}
\usepackage{enumitem}
\usepackage{titlesec}
\usepackage{fancyhdr}
\usepackage{xurl}
\usepackage{xcolor}
\definecolor{azul}{RGB}{25,70,140}
\usepackage[colorlinks=true, linkcolor=azul, urlcolor=azul, citecolor=azul, bookmarksnumbered=true,
  bookmarksopen=true, pdfdisplaydoctitle=true]{hyperref}
\urlstyle{same}

% Espaçamento simples; recuo de 1,25 cm na primeira linha (ABNT), sem espaço extra entre parágrafos
\setlength{\parindent}{1.25cm}
\setlength{\parskip}{0pt plus 1pt}
\setlength{\emergencystretch}{1.5em}
\frenchspacing  % espaço simples depois do ponto, como no português
% sem linhas órfãs (primeira linha do parágrafo no pé da página) nem viúvas (última no topo)
\clubpenalty=10000
\widowpenalty=10000
\displaywidowpenalty=10000
\setcounter{secnumdepth}{2}
\titleformat{\section}{\normalfont\normalsize\bfseries}{\thesection}{0.5em}{}
\titleformat{\subsection}{\normalfont\normalsize\bfseries}{\thesubsection}{0.5em}{}
\titlespacing*{\section}{0pt}{12pt}{0pt}
\titlespacing*{\subsection}{0pt}{12pt}{0pt}
\setlist[itemize]{leftmargin=1.2em, topsep=0pt, itemsep=1.5pt, parsep=0pt}
\providecommand{\tightlist}{\setlength{\itemsep}{1.5pt}\setlength{\parskip}{0pt}}

% Número da página no canto superior direito (ABNT)
\pagestyle{fancy}
\fancyhf{}
\fancyhead[R]{\footnotesize\thepage}
\renewcommand{\headrulewidth}{0pt}
\setlength{\headheight}{14.5pt}
\fancypagestyle{plain}{\fancyhf{}\fancyhead[R]{\footnotesize\thepage}}

% A Figura 1 ocupa uma página; floats grandes vão para a página seguinte à menção
\renewcommand{\floatpagefraction}{0.9}
\renewcommand{\topfraction}{0.86}
\renewcommand{\textfraction}{0.1}
\setlength{\LTpre}{6pt}
\setlength{\LTpost}{0pt}

% Referências (ABNT NBR 6023): alinhadas à esquerda, espaço simples, separadas por meia linha em branco (a linha
% inteira da norma levaria o artigo a 16 páginas, acima do limite de 15)
\newenvironment{referencias}{\raggedright\setlength{\parindent}{0pt}\setlength{\parskip}{0.5\baselineskip}}{\par}
\usepackage[section]{placeins}
\hypersetup{pdftitle={Declaração de uso de inteligência artificial}, pdfauthor={Felipe Carvalho Souza Santos}, pdfsubject={Avaliação de Políticas Públicas (PECO 5046-6046), PPGEco/UFES}, pdflang={pt-BR}}

\begin{document}

\thispagestyle{empty}
\begin{center}
{\bfseries UNIVERSIDADE FEDERAL DO ESPÍRITO SANTO\par}
{\bfseries CENTRO DE CIÊNCIAS JURÍDICAS E ECONÔMICAS\par}
{\bfseries PROGRAMA DE PÓS-GRADUAÇÃO EM ECONOMIA\par}
\vspace{12pt}
{\bfseries DECLARAÇÃO DE USO DE INTELIGÊNCIA ARTIFICIAL\par}
\end{center}
\vspace{6pt}
\begingroup\setlength{\parindent}{0pt}\setlength{\parskip}{6pt}\setlist[itemize]{itemsep=2pt, topsep=3pt}
Declaro, para os devidos fins, que ferramentas de Inteligência Artificial Generativa (IAG)
\textbf{\emph{foram}} utilizadas na elaboração deste trabalho nas seguintes etapas:

\begin{itemize}
\item
  A ferramenta Claude (Anthropic), operada pelo Claude Code com acesso ao repositório do trabalho,
  foi utilizada para ler e resumir o material da disciplina, as normas da LICC e a literatura, e
  para buscar e conferir as referências em bases bibliográficas (Crossref, OpenAlex) e páginas
  oficiais, na etapa de revisão da literatura e levantamento de fontes.
\item
  A ferramenta Claude (Anthropic) foi utilizada para escrever os programas de coleta, padronização e
  cruzamento dos dados públicos da LICC (anexos da SECULT, Portal da Transparência, Diário Oficial,
  Mapa Cultural e SALIC), inclusive uma automação de coleta no GitHub Actions e buscas com o serviço
  de raspagem Firecrawl, na etapa de coleta e tratamento dos dados.
\item
  A ferramenta Claude (Anthropic) foi utilizada para programar as estatísticas descritivas, as
  tabelas, a figura da teoria da mudança e o cálculo do poder estatístico, na etapa de análise dos
  dados.
\item
  A ferramenta Claude (Anthropic) foi utilizada para propor e discutir hipóteses e alternativas de
  desenho da avaliação de impacto, das quais o autor definiu o escopo e o desenho final, na etapa de
  desenho da avaliação.
\item
  A ferramenta Claude (Anthropic) foi utilizada para redigir versões preliminares das seções do
  artigo, que o autor editou, e para revisar gramática, concordância e estilo, na etapa de redação.
\item
  A ferramenta Claude (Anthropic), com os protocolos Academic Research Skills, foi utilizada para
  simular pareceres de revisão por pares e conferir cada número e referência do texto com as fontes,
  na etapa de revisão.
\item
  A ferramenta Claude (Anthropic) foi utilizada para gerar as versões em Word e LaTeX, a diagramação
  e os hiperlinks, na etapa de formatação.
\end{itemize}

Declaro ainda que todas as informações foram revisadas pelo autor. O autor assume integral
responsabilidade pela originalidade, veracidade e integridade do trabalho, inclusive por eventuais
imprecisões, inconsistências ou plágios provenientes do uso de tais ferramentas, em conformidade com
a Política de Integridade na Atividade Científica do CNPq estabelecida na Portaria CNPq n° 2.664, de
março de 2026.
\par\endgroup
\vspace{14pt}
\begin{center}
Vitória, 28 de setembro de 2026.\par
\textbf{Felipe Carvalho Souza Santos}\par
{\small Mini artigo ``Avaliação do desenho da Lei de Incentivo à Cultura Capixaba e proposta de avaliação de impacto''\par}
{\small Avaliação de Políticas Públicas (PECO 5046-6046), PPGEco/UFES\par}
\end{center}

\end{document}
